\documentclass[12pt,a4paper]{article}
\usepackage[utf8]{inputenc}
\usepackage[T1]{fontenc}
\usepackage{lmodern}
\usepackage[margin=1in]{geometry}
\usepackage{hyperref}
\usepackage{xcolor}
\usepackage{listings}
\usepackage{tcolorbox}
\usepackage{enumitem}
\usepackage{fancyhdr}
\usepackage{titlesec}
\usepackage{graphicx}
\usepackage{amssymb}
\usepackage{fontawesome5}

% Color definitions
\definecolor{codegreen}{rgb}{0,0.6,0}
\definecolor{codegray}{rgb}{0.5,0.5,0.5}
\definecolor{codepurple}{rgb}{0.58,0,0.82}
\definecolor{backcolour}{rgb}{0.97,0.97,0.97}
\definecolor{checkpoint}{rgb}{0.2,0.6,0.2}
\definecolor{gitcolor}{rgb}{0.9,0.4,0.1}
\definecolor{warningcolor}{rgb}{0.8,0.2,0.2}
\definecolor{conceptcolor}{rgb}{0.1,0.3,0.6}

% Code listing style
\lstdefinestyle{mystyle}{
    backgroundcolor=\color{backcolour},
    commentstyle=\color{codegreen},
    keywordstyle=\color{blue}\bfseries,
    numberstyle=\tiny\color{codegray},
    stringstyle=\color{codepurple},
    basicstyle=\ttfamily\small,
    breakatwhitespace=false,
    breaklines=true,
    captionpos=b,
    keepspaces=true,
    numbers=left,
    numbersep=5pt,
    showspaces=false,
    showstringspaces=false,
    showtabs=false,
    tabsize=2,
    frame=single,
    rulecolor=\color{codegray}
}
\lstset{style=mystyle}

% Custom environments
\newtcolorbox{checkpoint}{
    colback=green!5!white,
    colframe=checkpoint,
    title={\faCheckCircle\ \ CHECKPOINT: Verify Your Work},
    fonttitle=\bfseries
}

\newtcolorbox{gitcommit}{
    colback=orange!5!white,
    colframe=gitcolor,
    title={\faGitAlt\ \ GIT COMMIT REQUIRED},
    fonttitle=\bfseries
}

\newtcolorbox{concept}{
    colback=blue!5!white,
    colframe=conceptcolor,
    title={\faLightbulb\ \ FUNDAMENTAL CONCEPT},
    fonttitle=\bfseries
}

\newtcolorbox{warning}{
    colback=red!5!white,
    colframe=warningcolor,
    title={\faExclamationTriangle\ \ COMMON MISTAKE},
    fonttitle=\bfseries
}

\newtcolorbox{explanation}{
    colback=gray!5!white,
    colframe=codegray,
    title={\faCode\ \ LINE-BY-LINE EXPLANATION},
    fonttitle=\bfseries
}

% Header/Footer
\pagestyle{fancy}
\fancyhf{}
\rhead{React + TypeScript: First Principles}
\lhead{Building Your Personal Website}
\cfoot{\thepage}

\title{\Huge\textbf{Building a Personal Website from Scratch}\\[0.5cm]
\Large A First-Principles Guide to React and TypeScript\\[0.3cm]
\normalsize From Zero to Professional Portfolio}
\author{Comprehensive Tutorial for Beginners}
\date{January 2026}

\begin{document}

\maketitle
\thispagestyle{empty}

\vspace{1cm}

\begin{abstract}
This document provides an exhaustive, line-by-line guide to building a personal portfolio website using React and TypeScript. Written from the perspective of a computer science professor teaching fundamentals and a senior engineer ensuring production-quality code, every concept is explained from first principles. No prior programming knowledge is assumed. By the end, you will have built a professional website and deeply understand \textit{why} every piece of code exists.
\end{abstract}

\newpage
\tableofcontents
\newpage

%==============================================================================
% CHAPTER 1: FOUNDATIONS
%==============================================================================
\section{Chapter 1: What is Programming?}

Before we write a single line of code, we must understand what programming actually \textit{is}.

\begin{concept}
\textbf{Programming} is the act of giving precise instructions to a computer. A computer is fundamentally a machine that can:
\begin{enumerate}
    \item Store information (memory)
    \item Perform calculations (processing)
    \item Make decisions based on conditions (logic)
    \item Repeat actions (loops)
\end{enumerate}
A program is simply a sequence of these instructions written in a language the computer can understand.
\end{concept}

\subsection{Why Computers Need Special Languages}

Computers, at their core, only understand binary---sequences of 0s and 1s. Writing programs in binary would be impossibly tedious:

\begin{lstlisting}[language={},caption={Binary code (what computers actually understand)}]
01001000 01100101 01101100 01101100 01101111
\end{lstlisting}

This is why we invented \textbf{programming languages}---human-readable ways to express instructions that get translated into binary.

\subsection{The Language Hierarchy}

\begin{enumerate}
    \item \textbf{Machine Code}: Raw binary (0s and 1s)
    \item \textbf{Assembly}: Human-readable but still very low-level
    \item \textbf{High-Level Languages}: JavaScript, Python, TypeScript---what we use
\end{enumerate}

\begin{concept}
\textbf{JavaScript} is the programming language that web browsers understand natively. When you visit any website, JavaScript is what makes it interactive---handling button clicks, animations, and data updates.

\textbf{TypeScript} is JavaScript with an added safety layer called ``types.'' We'll explain this deeply later.
\end{concept}

\subsection{What is a Website?}

A website consists of three fundamental technologies:

\begin{enumerate}
    \item \textbf{HTML (HyperText Markup Language)}: The structure/content
    \item \textbf{CSS (Cascading Style Sheets)}: The visual appearance
    \item \textbf{JavaScript}: The behavior/interactivity
\end{enumerate}

Think of building a house:
\begin{itemize}
    \item HTML = The walls, rooms, and foundation (structure)
    \item CSS = The paint, furniture, and decorations (appearance)
    \item JavaScript = The electrical system, plumbing, and automation (functionality)
\end{itemize}

%==============================================================================
% CHAPTER 2: SETTING UP
%==============================================================================
\section{Chapter 2: Setting Up Your Development Environment}

Before cooking, you need a kitchen. Before coding, you need a development environment.

\subsection{Step 2.1: Install Node.js}

\begin{concept}
\textbf{Node.js} is a ``runtime environment'' for JavaScript. Normally, JavaScript only runs inside web browsers. Node.js allows JavaScript to run on your computer directly, which we need for development tools.

Think of it this way: A browser is like a fish tank where JavaScript fish can swim. Node.js builds another tank on your computer so JavaScript fish can swim there too.
\end{concept}

\textbf{Action Required:}
\begin{enumerate}
    \item Open your web browser
    \item Go to: \url{https://nodejs.org}
    \item Download the \textbf{LTS} (Long Term Support) version
    \item Run the installer and accept all defaults
\end{enumerate}

\begin{checkpoint}
Open your Terminal (Mac) or Command Prompt (Windows) and type:
\begin{lstlisting}[language=bash]
node --version
\end{lstlisting}
You should see something like: \texttt{v20.x.x} or higher.

Then type:
\begin{lstlisting}[language=bash]
npm --version
\end{lstlisting}
You should see a version number like \texttt{10.x.x}.

If both commands show version numbers, Node.js is installed correctly. If you see ``command not found,'' restart your terminal and try again.
\end{checkpoint}

\begin{concept}
\textbf{npm} (Node Package Manager) comes bundled with Node.js. It's a tool that lets you download and use code that other programmers have written. Instead of building everything from scratch, you can use ``packages''---pre-written solutions to common problems.

This is like using pre-made ingredients when cooking instead of growing all your vegetables yourself.
\end{concept}

\subsection{Step 2.2: Install Visual Studio Code}

\textbf{Visual Studio Code (VS Code)} is a text editor designed for programming. While you \textit{could} write code in Notepad, VS Code provides:
\begin{itemize}
    \item Syntax highlighting (colors different parts of code)
    \item Error detection (underlines mistakes)
    \item Auto-completion (suggests code as you type)
    \item Integrated terminal (run commands without leaving the editor)
\end{itemize}

\textbf{Action Required:}
\begin{enumerate}
    \item Go to: \url{https://code.visualstudio.com}
    \item Download and install for your operating system
    \item Open VS Code after installation
\end{enumerate}

\subsection{Step 2.3: Install Essential VS Code Extensions}

Extensions add superpowers to VS Code.

\textbf{Action Required:}
\begin{enumerate}
    \item In VS Code, click the Extensions icon (four squares) on the left sidebar
    \item Search for and install these extensions:
    \begin{itemize}
        \item \textbf{ESLint} - catches coding mistakes
        \item \textbf{Prettier} - formats your code consistently
        \item \textbf{TypeScript and JavaScript Language Features} - usually pre-installed
    \end{itemize}
\end{enumerate}

\begin{checkpoint}
You should now have:
\begin{itemize}
    \item[$\square$] Node.js installed (verified with \texttt{node --version})
    \item[$\square$] npm installed (verified with \texttt{npm --version})
    \item[$\square$] VS Code installed and open
    \item[$\square$] ESLint and Prettier extensions installed
\end{itemize}
\end{checkpoint}

\subsection{Step 2.4: Install Git}

\begin{concept}
\textbf{Git} is a ``version control system.'' It tracks every change you make to your code, allowing you to:
\begin{itemize}
    \item Go back to any previous version if something breaks
    \item See exactly what changed and when
    \item Collaborate with others without overwriting each other's work
\end{itemize}

Think of Git as an ``undo history'' for your entire project that never forgets.
\end{concept}

\textbf{Action Required:}
\begin{enumerate}
    \item Go to: \url{https://git-scm.com/downloads}
    \item Download and install for your operating system
    \item Accept all default options during installation
\end{enumerate}

\begin{checkpoint}
Open a new terminal and type:
\begin{lstlisting}[language=bash]
git --version
\end{lstlisting}
You should see something like: \texttt{git version 2.x.x}
\end{checkpoint}

\subsection{Step 2.5: Create a GitHub Account}

\begin{concept}
\textbf{GitHub} is a website that hosts Git repositories (projects) online. It allows you to:
\begin{itemize}
    \item Back up your code to the cloud
    \item Share your work publicly (like a portfolio)
    \item Collaborate with others
\end{itemize}
Git is the tool; GitHub is the online storage.
\end{concept}

\textbf{Action Required:}
\begin{enumerate}
    \item Go to: \url{https://github.com}
    \item Click ``Sign Up'' and create an account
    \item Verify your email address
\end{enumerate}

%==============================================================================
% CHAPTER 3: CREATING THE PROJECT
%==============================================================================
\section{Chapter 3: Creating Your React Project}

Now we create the actual project. We'll use a tool called \textbf{Vite} (pronounced ``veet'').

\begin{concept}
\textbf{Vite} is a ``build tool'' that:
\begin{enumerate}
    \item Creates a project structure for you
    \item Runs a local development server (so you can see your website)
    \item Bundles your code for production (prepares it for the real internet)
\end{enumerate}
It's like a sous chef that handles all the prep work so you can focus on cooking.
\end{concept}

\subsection{Step 3.1: Choose Your Project Location}

\textbf{Action Required:}
\begin{enumerate}
    \item Open Terminal (Mac) or Command Prompt (Windows)
    \item Navigate to where you want your project:
\end{enumerate}

\begin{lstlisting}[language=bash,caption={Navigate to your desired folder}]
cd ~/Documents
\end{lstlisting}

\begin{explanation}
\textbf{Line 1:} \texttt{cd} stands for ``change directory.'' The \texttt{\~} symbol means ``home folder.'' So this command moves you to your Documents folder.
\end{explanation}

\subsection{Step 3.2: Create the Project}

\begin{lstlisting}[language=bash,caption={Create a new React + TypeScript project}]
npm create vite@latest my-portfolio -- --template react-ts
\end{lstlisting}

\begin{explanation}
Let's break down every part of this command:

\begin{itemize}
    \item \texttt{npm} - Use the Node Package Manager
    \item \texttt{create} - npm's command to run project generators
    \item \texttt{vite@latest} - Use the latest version of Vite
    \item \texttt{my-portfolio} - The name of your project folder (you can change this)
    \item \texttt{--} - Separates npm arguments from Vite arguments
    \item \texttt{--template react-ts} - Use the React + TypeScript template
\end{itemize}
\end{explanation}

\subsection{Step 3.3: Enter the Project and Install Dependencies}

\begin{lstlisting}[language=bash,caption={Navigate into project and install}]
cd my-portfolio
npm install
\end{lstlisting}

\begin{explanation}
\textbf{Line 1:} Move into the newly created project folder.

\textbf{Line 2:} \texttt{npm install} reads a file called \texttt{package.json} (which Vite created) and downloads all the packages your project needs. This creates a \texttt{node\_modules} folder containing thousands of files---this is normal!
\end{explanation}

\begin{warning}
The \texttt{node\_modules} folder will be huge (hundreds of MB). This is expected. Never manually edit files inside \texttt{node\_modules}---they're automatically managed.
\end{warning}

\subsection{Step 3.4: Start the Development Server}

\begin{lstlisting}[language=bash,caption={Start the local development server}]
npm run dev
\end{lstlisting}

\begin{explanation}
This command:
\begin{enumerate}
    \item Starts a local web server on your computer
    \item Watches for changes to your files
    \item Automatically refreshes the browser when you save changes
\end{enumerate}
\end{explanation}

\begin{checkpoint}
After running \texttt{npm run dev}, you should see output like:
\begin{lstlisting}[language=bash]
  VITE v5.x.x  ready in xxx ms

  -> Local:   http://localhost:5173/
  -> Network: use --host to expose
\end{lstlisting}

\textbf{Open your web browser} and go to \texttt{http://localhost:5173/}

You should see a page with the Vite and React logos, and a counter button. Click the button---it should increment. This proves React is working!
\end{checkpoint}

\begin{concept}
\textbf{localhost} means ``this computer.'' \textbf{5173} is a ``port number''---think of it as a specific door into your computer. The development server opens door 5173 and serves your website through it.
\end{concept}

\subsection{Step 3.5: Open in VS Code}

Keep the terminal running (the server needs to stay on), and open a \textbf{new} terminal window/tab.

\begin{lstlisting}[language=bash,caption={Open project in VS Code}]
cd ~/Documents/my-portfolio
code .
\end{lstlisting}

\begin{explanation}
\textbf{Line 1:} Navigate to your project folder.

\textbf{Line 2:} \texttt{code .} opens VS Code in the current directory. The \texttt{.} means ``here.''
\end{explanation}

\begin{checkpoint}
You should now have:
\begin{itemize}
    \item[$\square$] Terminal with server running (showing localhost URL)
    \item[$\square$] Browser showing the Vite + React default page
    \item[$\square$] VS Code open with your project files visible in the left sidebar
\end{itemize}
\end{checkpoint}

%==============================================================================
% CHAPTER 4: UNDERSTANDING THE PROJECT STRUCTURE
%==============================================================================
\section{Chapter 4: Understanding the Project Structure}

Before modifying anything, let's understand what Vite created.

\subsection{The File Tree}

Your project should look like this:
\begin{lstlisting}[language=bash,caption={Project structure}]
my-portfolio/
|-- node_modules/      # Downloaded packages (don't touch)
|-- public/            # Static files served as-is
|-- src/               # YOUR CODE GOES HERE
|   |-- assets/        # Images, fonts, etc.
|   |-- App.css        # Styles for App component
|   |-- App.tsx        # Main App component
|   |-- index.css      # Global styles
|   |-- main.tsx       # Entry point
|   |-- vite-env.d.ts  # TypeScript definitions for Vite
|-- .gitignore         # Files Git should ignore
|-- index.html         # The HTML shell
|-- package.json       # Project configuration
|-- tsconfig.json      # TypeScript configuration
|-- vite.config.ts     # Vite configuration
\end{lstlisting}

\subsection{Understanding Each File}

\subsubsection{index.html - The HTML Shell}

Open \texttt{index.html} in VS Code:

\begin{lstlisting}[language=html,caption={index.html}]
<!DOCTYPE html>
<html lang="en">
  <head>
    <meta charset="UTF-8" />
    <link rel="icon" type="image/svg+xml" href="/vite.svg" />
    <meta name="viewport" content="width=device-width, initial-scale=1.0" />
    <title>Vite + React + TS</title>
  </head>
  <body>
    <div id="root"></div>
    <script type="module" src="/src/main.tsx"></script>
  </body>
</html>
\end{lstlisting}

\begin{explanation}
\textbf{Line 1:} \texttt{<!DOCTYPE html>} tells browsers this is an HTML5 document.

\textbf{Line 2:} \texttt{<html lang="en">} is the root element. \texttt{lang="en"} tells browsers/screen readers the content is in English.

\textbf{Lines 3-8:} The \texttt{<head>} contains metadata (information \textit{about} the page, not visible content):
\begin{itemize}
    \item \texttt{charset="UTF-8"} - Use Unicode character encoding (supports all languages/emojis)
    \item \texttt{<link rel="icon">} - The favicon (tab icon)
    \item \texttt{viewport} meta tag - Makes the site responsive on mobile
    \item \texttt{<title>} - Text shown in the browser tab
\end{itemize}

\textbf{Line 10:} \texttt{<div id="root">} - This empty div is crucial! React will insert your entire application inside this element.

\textbf{Line 11:} \texttt{<script type="module" src="/src/main.tsx">} - Loads your React application. \texttt{type="module"} enables modern JavaScript features.
\end{explanation}

\subsubsection{src/main.tsx - The Entry Point}

\begin{lstlisting}[language=JavaScript,caption={src/main.tsx}]
import React from 'react'
import ReactDOM from 'react-dom/client'
import App from './App.tsx'
import './index.css'

ReactDOM.createRoot(document.getElementById('root')!).render(
  <React.StrictMode>
    <App />
  </React.StrictMode>,
)
\end{lstlisting}

\begin{explanation}
\textbf{Line 1:} \texttt{import React from 'react'} - Import the React library. \texttt{import} is how we bring in code from other files or packages.

\textbf{Line 2:} \texttt{import ReactDOM from 'react-dom/client'} - Import ReactDOM, which connects React to the actual web page (the DOM - Document Object Model).

\textbf{Line 3:} \texttt{import App from './App.tsx'} - Import our main App component. \texttt{./} means ``in the current folder.''

\textbf{Line 4:} \texttt{import './index.css'} - Import global CSS styles.

\textbf{Line 6:} This is where the magic happens:
\begin{itemize}
    \item \texttt{document.getElementById('root')} - Find the div with id="root" in index.html
    \item \texttt{!} - TypeScript syntax meaning ``I promise this exists''
    \item \texttt{ReactDOM.createRoot(...)} - Create a React rendering root at that element
    \item \texttt{.render(...)} - Actually render our React components there
\end{itemize}

\textbf{Lines 7-9:} \texttt{<React.StrictMode>} wraps our app and enables extra development warnings. \texttt{<App />} is our main component.
\end{explanation}

\begin{concept}
\textbf{What is a Component?}

A React component is a reusable piece of UI (user interface). Think of components like LEGO bricks:
\begin{itemize}
    \item Each brick (component) is self-contained
    \item You can combine bricks to build complex structures
    \item You can reuse the same brick multiple times
\end{itemize}

Your website will be made of components: Header component, About component, Experience component, etc.
\end{concept}

\subsubsection{src/App.tsx - Your First Component}

\begin{lstlisting}[language=JavaScript,caption={src/App.tsx (default)}]
import { useState } from 'react'
import reactLogo from './assets/react.svg'
import viteLogo from '/vite.svg'
import './App.css'

function App() {
  const [count, setCount] = useState(0)

  return (
    <>
      <div>
        <a href="https://vite.dev" target="_blank">
          <img src={viteLogo} className="logo" alt="Vite logo" />
        </a>
        <a href="https://react.dev" target="_blank">
          <img src={reactLogo} className="logo react" alt="React logo" />
        </a>
      </div>
      <h1>Vite + React</h1>
      <div className="card">
        <button onClick={() => setCount((count) => count + 1)}>
          count is {count}
        </button>
        <p>
          Edit <code>src/App.tsx</code> and save to test HMR
        </p>
      </div>
      <p className="read-the-docs">
        Click on the Vite and React logos to learn more
      </p>
    </>
  )
}

export default App
\end{lstlisting}

\begin{explanation}
\textbf{Line 1:} \texttt{import \{ useState \} from 'react'} - Import the \texttt{useState} function from React. The curly braces mean we're importing a \textit{specific} export, not the default export.

\textbf{Lines 2-4:} Import images and CSS.

\textbf{Line 6:} \texttt{function App()} - Define a function called App. In React, components are functions that return what should be displayed.

\textbf{Line 7:} \texttt{const [count, setCount] = useState(0)} - This is a ``state hook'':
\begin{itemize}
    \item \texttt{count} - A variable holding the current value (starts at 0)
    \item \texttt{setCount} - A function to update that value
    \item \texttt{useState(0)} - Initialize state with value 0
\end{itemize}

\textbf{Lines 9-31:} The \texttt{return} statement contains JSX---HTML-like syntax inside JavaScript.

\textbf{Line 10:} \texttt{<>...</>} is a ``Fragment''---a way to group elements without adding extra HTML.

\textbf{Line 21:} \texttt{onClick=\{() => setCount((count) => count + 1)\}} - When clicked, increase count by 1.

\textbf{Line 34:} \texttt{export default App} - Make this component available to other files. This is why main.tsx can \texttt{import App from './App.tsx'}.
\end{explanation}

\begin{concept}
\textbf{What is JSX?}

JSX looks like HTML but it's actually JavaScript in disguise. When your code compiles, this:
\begin{lstlisting}[language=JavaScript]
<h1>Hello</h1>
\end{lstlisting}
Becomes this:
\begin{lstlisting}[language=JavaScript]
React.createElement('h1', null, 'Hello')
\end{lstlisting}
JSX is syntactic sugar---it makes React code more readable.
\end{concept}

\begin{concept}
\textbf{What is State?}

State is data that can change over time. When state changes, React automatically re-renders (redraws) the component to show the new data.

The counter button demonstrates this:
\begin{enumerate}
    \item Initial state: count = 0
    \item User clicks button
    \item setCount is called, count becomes 1
    \item React sees state changed, re-renders
    \item Screen now shows ``count is 1''
\end{enumerate}
\end{concept}

%==============================================================================
% CHAPTER 5: GIT INITIALIZATION
%==============================================================================
\section{Chapter 5: Initialize Git Repository}

Before making changes, let's set up version control.

\subsection{Step 5.1: Initialize Git}

Open the VS Code integrated terminal (View → Terminal or \texttt{Ctrl+`}):

\begin{lstlisting}[language=bash,caption={Initialize Git repository}]
git init
\end{lstlisting}

\begin{explanation}
\texttt{git init} creates a hidden \texttt{.git} folder that stores all version control data. Your project is now a Git repository.
\end{explanation}

\subsection{Step 5.2: Configure Git (First Time Only)}

\begin{lstlisting}[language=bash,caption={Set your identity}]
git config --global user.name "Your Name"
git config --global user.email "your.email@example.com"
\end{lstlisting}

\begin{explanation}
Git needs to know who you are so it can record who made each change. Replace with your actual name and email (use the email associated with your GitHub account).
\end{explanation}

\subsection{Step 5.3: Understand .gitignore}

Open \texttt{.gitignore}:

\begin{lstlisting}[caption={.gitignore (partial)}]
# Logs
logs
*.log
npm-debug.log*

# Dependencies
node_modules

# Build output
dist

# Local env files
.env.local
.env.*.local
\end{lstlisting}

\begin{explanation}
\texttt{.gitignore} lists files and folders that Git should NOT track:
\begin{itemize}
    \item \texttt{node\_modules} - Too large, can be regenerated with \texttt{npm install}
    \item \texttt{dist} - Build output, regenerated when you build
    \item \texttt{.env} files - Often contain secrets (passwords, API keys)
\end{itemize}
\end{explanation}

\subsection{Step 5.4: Make Your First Commit}

\begin{lstlisting}[language=bash,caption={Stage and commit all files}]
git add .
git commit -m "Initial commit: Vite + React + TypeScript setup"
\end{lstlisting}

\begin{explanation}
\textbf{Line 1:} \texttt{git add .} - Stage all files for commit. The \texttt{.} means ``everything in current directory.'' Staging is like putting items in a shopping cart before checkout.

\textbf{Line 2:} \texttt{git commit -m "message"} - Create a snapshot of the staged files. The \texttt{-m} flag lets you include a message describing what changed.
\end{explanation}

\begin{gitcommit}
You've made your first commit! Verify with:
\begin{lstlisting}[language=bash]
git log --oneline
\end{lstlisting}
You should see your commit with its message.
\end{gitcommit}

\subsection{Step 5.5: Connect to GitHub}

\begin{enumerate}
    \item Go to \url{https://github.com}
    \item Click the ``+'' icon → ``New repository''
    \item Name it \texttt{my-portfolio} (or your chosen name)
    \item Keep it public (for portfolio visibility)
    \item Do NOT initialize with README (we already have files)
    \item Click ``Create repository''
\end{enumerate}

GitHub will show commands. Run these (replace USERNAME with your GitHub username):

\begin{lstlisting}[language=bash,caption={Connect local repo to GitHub}]
git remote add origin https://github.com/USERNAME/my-portfolio.git
git branch -M main
git push -u origin main
\end{lstlisting}

\begin{explanation}
\textbf{Line 1:} Connect your local repository to the GitHub repository. ``origin'' is the conventional name for the main remote.

\textbf{Line 2:} Rename the default branch to ``main'' (Git's new standard).

\textbf{Line 3:} Push (upload) your code to GitHub. \texttt{-u} sets up tracking so future pushes are simpler.
\end{explanation}

\begin{checkpoint}
\begin{itemize}
    \item[$\square$] Refresh your GitHub repository page---you should see your files
    \item[$\square$] The commit message ``Initial commit'' should appear
\end{itemize}
\end{checkpoint}

%==============================================================================
% CHAPTER 6: TYPESCRIPT FUNDAMENTALS
%==============================================================================
\section{Chapter 6: TypeScript Fundamentals}

Before building components, you must understand TypeScript.

\begin{concept}
\textbf{What is TypeScript?}

JavaScript is ``dynamically typed''---variables can hold any type of data:
\begin{lstlisting}[language=JavaScript]
let x = 5;        // x is a number
x = "hello";      // now x is a string - JavaScript allows this!
\end{lstlisting}

TypeScript adds ``static types''---you declare what type a variable holds:
\begin{lstlisting}[language=JavaScript]
let x: number = 5;   // x must be a number
x = "hello";         // ERROR! TypeScript catches this mistake
\end{lstlisting}

This catches bugs before your code even runs.
\end{concept}

\subsection{Basic Types}

\begin{lstlisting}[language=JavaScript,caption={TypeScript basic types}]
// Primitive types
let name: string = "Sofia";
let age: number = 17;
let isStudent: boolean = true;

// Arrays
let skills: string[] = ["React", "TypeScript", "Python"];
let numbers: number[] = [1, 2, 3, 4, 5];

// Objects
let person: { name: string; age: number } = {
  name: "Sofia",
  age: 17
};
\end{lstlisting}

\begin{explanation}
\textbf{Lines 2-4:} Basic types - string (text), number (any number), boolean (true/false).

\textbf{Lines 7-8:} Arrays use \texttt{type[]} syntax. \texttt{string[]} means ``array of strings.''

\textbf{Lines 11-14:} Objects define their structure with curly braces, listing each property and its type.
\end{explanation}

\subsection{Type Inference}

\begin{lstlisting}[language=JavaScript,caption={TypeScript infers types automatically}]
// You don't always need to specify types
let message = "Hello";  // TypeScript knows this is a string
let count = 42;         // TypeScript knows this is a number

// But being explicit is clearer for complex types
let user: { name: string; email: string } = {
  name: "Sofia",
  email: "sofia@example.com"
};
\end{lstlisting}

\begin{concept}
TypeScript is smart---it can \textit{infer} (figure out) types from the values you assign. But explicit types serve as documentation and catch errors earlier.
\end{concept}

\subsection{Interfaces}

\begin{lstlisting}[language=JavaScript,caption={Defining reusable type shapes}]
// Interface defines the shape of an object
interface Experience {
  title: string;
  company: string;
  startDate: string;
  endDate?: string;      // ? means optional
  description: string;
}

// Now we can use this interface
const job: Experience = {
  title: "Machine Learning Intern",
  company: "Convictional",
  startDate: "Jul 2025",
  endDate: "Aug 2025",
  description: "Researched multi-vector retrieval methods"
};
\end{lstlisting}

\begin{explanation}
\textbf{Lines 2-8:} \texttt{interface} creates a reusable type definition. Any object of type \texttt{Experience} must have these properties.

\textbf{Line 6:} The \texttt{?} makes a property optional. \texttt{endDate} can be omitted (useful for current jobs).

\textbf{Lines 11-17:} Using the interface ensures our object has the correct shape.
\end{explanation}

\subsection{Functions with Types}

\begin{lstlisting}[language=JavaScript,caption={Typed function parameters and return values}]
// Function with typed parameters and return type
function greet(name: string): string {
  return `Hello, ${name}!`;
}

// Arrow function syntax (common in React)
const add = (a: number, b: number): number => {
  return a + b;
};

// Function that doesn't return anything
const logMessage = (message: string): void => {
  console.log(message);
};
\end{lstlisting}

\begin{explanation}
\textbf{Line 2:} \texttt{(name: string)} - parameter must be a string. \texttt{: string} after parentheses - function returns a string.

\textbf{Line 7:} Arrow functions are a shorter syntax. \texttt{=>} is the ``arrow.''

\textbf{Line 12:} \texttt{void} means the function returns nothing.
\end{explanation}

%==============================================================================
% CHAPTER 7: BUILDING YOUR FIRST COMPONENT
%==============================================================================
\section{Chapter 7: Building Your First Real Component}

Now let's clear out the default content and start building YOUR website.

\subsection{Step 7.1: Clean Up App.tsx}

Replace the entire contents of \texttt{src/App.tsx}:

\begin{lstlisting}[language=JavaScript,caption={src/App.tsx - Clean starting point}]
function App() {
  return (
    <main>
      <h1>sofia bodnar</h1>
      <p>17 y/o high school student</p>
    </main>
  )
}

export default App
\end{lstlisting}

\begin{explanation}
\textbf{Line 1:} Define the App component as a function.

\textbf{Lines 2-7:} Return JSX. We use \texttt{<main>} as the container element (semantic HTML for the main content area).

\textbf{Line 4-5:} Your name and tagline as HTML elements.

\textbf{Line 10:} Export the component so other files can import it.
\end{explanation}

\begin{checkpoint}
Save the file. Your browser should automatically refresh and show just your name and tagline on a plain white page.

If it doesn't refresh, check that \texttt{npm run dev} is still running in your terminal.
\end{checkpoint}

\subsection{Step 7.2: Clean Up CSS}

Replace \texttt{src/index.css}:

\begin{lstlisting}[language=CSS,caption={src/index.css - Clean global styles}]
* {
  margin: 0;
  padding: 0;
  box-sizing: border-box;
}

body {
  font-family: -apple-system, BlinkMacSystemFont, 'Segoe UI', Roboto,
    Oxygen, Ubuntu, Cantarell, sans-serif;
  line-height: 1.6;
  color: #333;
  background-color: #fafafa;
}
\end{lstlisting}

\begin{explanation}
\textbf{Lines 1-5:} The ``CSS reset.'' By default, browsers add margins and padding to elements. This removes them so we have full control.
\begin{itemize}
    \item \texttt{*} selects ALL elements
    \item \texttt{box-sizing: border-box} makes sizing calculations more intuitive
\end{itemize}

\textbf{Lines 7-13:} Base styles for the body:
\begin{itemize}
    \item \texttt{font-family} - System fonts that look native on each OS
    \item \texttt{line-height: 1.6} - Space between lines (1.6 times the font size)
    \item \texttt{color: \#333} - Dark gray text (softer than pure black)
    \item \texttt{background-color: \#fafafa} - Very light gray background
\end{itemize}
\end{explanation}

Replace \texttt{src/App.css} with an empty file or delete all its contents (we'll add styles later).

\begin{checkpoint}
Your page should now show:
\begin{itemize}
    \item Your name in bold (default h1 style)
    \item Tagline below it
    \item Light gray background
    \item No weird margins or spacing
\end{itemize}
\end{checkpoint}

\begin{gitcommit}
Good stopping point! Let's commit:
\begin{lstlisting}[language=bash]
git add .
git commit -m "Clean slate: Remove default Vite content, add basic structure"
git push
\end{lstlisting}
\end{gitcommit}

%==============================================================================
% CHAPTER 8: PROJECT ARCHITECTURE
%==============================================================================
\section{Chapter 8: Planning Your Component Architecture}

Before writing more code, let's plan our component structure.

\begin{concept}
\textbf{Component Thinking}

Look at your design. We need to identify:
\begin{enumerate}
    \item What pieces repeat? (Those become reusable components)
    \item What pieces are distinct sections? (Those become section components)
    \item What data does each piece need? (Those become props)
\end{enumerate}
\end{concept}

\subsection{Analyzing Your Design}

Looking at your mockup, I identify these components:

\begin{enumerate}
    \item \textbf{Header} - Name, tagline, beliefs
    \item \textbf{ThemeToggle} - Dark/light mode switch
    \item \textbf{SkillsGraph} - The visualization of your skills
    \item \textbf{Section} - Reusable section container (work, research, other)
    \item \textbf{ExperienceCard} - Each work experience entry
    \item \textbf{ResearchCard} - Each research entry
    \item \textbf{ProjectCard} - Each project entry
    \item \textbf{SocialLinks} - Footer with icons
\end{enumerate}

\subsection{Create the Folder Structure}

\begin{lstlisting}[language=bash,caption={Create component folders}]
mkdir -p src/components
mkdir -p src/data
mkdir -p src/styles
mkdir -p src/assets/icons
\end{lstlisting}

\begin{explanation}
\textbf{Line 1:} Create a components folder. \texttt{-p} means ``create parent directories if needed.''

\textbf{Line 2:} Data folder will hold your content (experiences, projects, etc.).

\textbf{Line 3:} Styles folder for CSS files.

\textbf{Line 4:} Assets/icons for SVG icons.
\end{explanation}

%==============================================================================
% CHAPTER 9: DATA-DRIVEN DESIGN
%==============================================================================
\section{Chapter 9: Data-Driven Design}

Professional websites separate \textit{data} from \textit{presentation}. This makes updates easy---change the data file, and the whole site updates.

\subsection{Step 9.1: Create Your Data File}

Create \texttt{src/data/content.ts}:

\begin{lstlisting}[language=JavaScript,caption={src/data/content.ts}]
// Define the shape of our data with TypeScript interfaces

export interface Experience {
  title: string;
  company: string;
  period: string;
  description: string;
  tags?: string[];
  logo?: string;
}

export interface Research {
  title: string;
  venue?: string;
  note?: string;
  description?: string;
}

export interface Project {
  title: string;
  tagline: string;
  metrics?: string;
}

// Your actual data
export const personalInfo = {
  name: "sofia bodnar",
  tagline: "17 y/o high school student",
  obsessions: ["machine learning", "coding", "math", "digits of pi"],
  belief: {
    text: "we are entering the",
    link: { text: "era of experience", url: "#" }
  }
};

export const skills = [
  "continual learning",
  "temporal knowledge graphs",
  "reinforcement learning",
  "GP processes",
  "test time training",
  "Darwin Godel Machine"
];

export const experiences: Experience[] = [
  {
    title: "Machine Learning Research Intern",
    company: "Sentra",
    period: "Oct 2025 - Current",
    description: "building a self improving memory architecture that learns from experience and optimizes curation policies via online Bayesian Optimization",
    tags: ["A16Z", "MIT"]
  },
  {
    title: "Software Engineer",
    company: "Omen",
    period: "Sep 2025 - Current",
    description: "developing end to end automation layer for incorporating startups and cap table management system"
  },
  {
    title: "Machine Learning Intern",
    company: "Convictional",
    period: "Jul 2025 - Aug 2025",
    description: "researched multi vector retrieval methods such as MUVERA and Colbert v2, analyzing trade-offs in latency, memory usage, and retrieval accuracy.",
    tags: ["Y Combinator"]
  }
];

export const research: Research[] = [
  {
    title: "Learning from Experience: Rethinking Information for Computationally Bounded Intelligence",
    venue: "ICML in Seoul, South Korea",
    note: "presented my research with MIT prof"
  },
  {
    title: "Spatial Clustering and Changepoint Detection in Forecast Failure Analysis",
    venue: "Canadian Science Publishing",
    note: "won monetary award",
    description: "using spatial autocorrelation (Moran's I, LISA), time series analysis (CUSUM), and tree-based models to detect systematic forecast errors before crises escalate in vulnerable regions."
  }
];

export const projects: Project[] = [
  {
    title: "Filey",
    tagline: "The first legal paperwork infrastructure for startups in Canada",
    metrics: "15k users, $25k MRR"
  },
  {
    title: "Axiom Competition",
    tagline: "helped cracked youth break into startup tech scene"
  }
];

export const socialLinks = {
  github: "https://github.com/yourusername",
  twitter: "https://twitter.com/yourusername",
  linkedin: "https://linkedin.com/in/yourusername",
  email: "mailto:your.email@example.com"
};
\end{lstlisting}

\begin{explanation}
\textbf{Lines 3-11:} Define the \texttt{Experience} interface. Every experience entry must have title, company, period, description. Tags and logo are optional (\texttt{?}).

\textbf{Lines 13-17:} Research interface.

\textbf{Lines 19-22:} Project interface.

\textbf{Lines 25-31:} Your personal information as a JavaScript object.

\textbf{Lines 33-38:} Skills array---simple list of strings.

\textbf{Lines 40-62:} Experiences array. Note \texttt{Experience[]}---this means ``array of Experience objects.'' TypeScript will error if any object doesn't match the interface.

\textbf{Lines 64-77:} Research entries.

\textbf{Lines 79-88:} Projects.

\textbf{Lines 90-95:} Social media links.
\end{explanation}

\begin{concept}
\textbf{Why Separate Data from Components?}

\begin{enumerate}
    \item \textbf{Single Source of Truth}: Update data in one place, changes reflect everywhere
    \item \textbf{Type Safety}: TypeScript ensures data has correct shape
    \item \textbf{Testability}: Data can be tested independently
    \item \textbf{Scalability}: Easy to load data from an API later
\end{enumerate}
\end{concept}

\begin{gitcommit}
\begin{lstlisting}[language=bash]
git add .
git commit -m "Add data layer with TypeScript interfaces and content"
git push
\end{lstlisting}
\end{gitcommit}

%==============================================================================
% CHAPTER 10: BUILDING THE HEADER
%==============================================================================
\section{Chapter 10: Building the Header Component}

\subsection{Step 10.1: Create the Header Component}

Create \texttt{src/components/Header.tsx}:

\begin{lstlisting}[language=JavaScript,caption={src/components/Header.tsx}]
import { personalInfo } from '../data/content';

interface HeaderProps {
  // We'll add props here later if needed
}

function Header({}: HeaderProps) {
  const { name, tagline, obsessions, belief } = personalInfo;

  return (
    <header className="header">
      <h1 className="header__name">{name}</h1>
      <p className="header__tagline">{tagline}</p>

      <p className="header__obsessions">
        obsessed with: {obsessions.join(', ')}
      </p>

      <p className="header__belief">
        what i strongly believe in: {belief.text}{' '}
        <a href={belief.link.url} className="header__link">
          {belief.link.text}
        </a>
      </p>
    </header>
  );
}

export default Header;
\end{lstlisting}

\begin{explanation}
\textbf{Line 1:} Import your data. The path \texttt{../data/content} means ``go up one folder, then into data, then content.''

\textbf{Lines 3-5:} Define a TypeScript interface for props (properties passed to the component). Currently empty but good practice to have.

\textbf{Line 7:} \texttt{function Header(\{\}: HeaderProps)} - Define component. \texttt{\{\}} destructures the props object (empty for now).

\textbf{Line 8:} \texttt{const \{ name, tagline, obsessions, belief \} = personalInfo;} - This is ``destructuring.'' It extracts properties from the object into individual variables.

Equivalent to:
\begin{lstlisting}[language=JavaScript]
const name = personalInfo.name;
const tagline = personalInfo.tagline;
// etc.
\end{lstlisting}

\textbf{Line 12:} \texttt{\{name\}} in JSX - Curly braces insert JavaScript values into the HTML.

\textbf{Line 16:} \texttt{obsessions.join(', ')} - The \texttt{join} method combines array elements into a string with the specified separator.

\textbf{Line 20:} \texttt{\{' '\}} - Explicit space in JSX (regular spaces between elements are sometimes ignored).

\textbf{Lines 21-23:} The anchor tag with href from our data.
\end{explanation}

\begin{concept}
\textbf{BEM Naming Convention}

Notice the CSS class names: \texttt{header}, \texttt{header\_\_name}, \texttt{header\_\_tagline}

This is \textbf{BEM} (Block Element Modifier):
\begin{itemize}
    \item \textbf{Block}: The component itself (\texttt{header})
    \item \textbf{Element}: Parts of the block (\texttt{header\_\_name})
    \item \textbf{Modifier}: Variations (\texttt{header\_\_name--large})
\end{itemize}

BEM keeps CSS organized and prevents naming conflicts.
\end{concept}

\subsection{Step 10.2: Create Header Styles}

Create \texttt{src/styles/Header.css}:

\begin{lstlisting}[language=CSS,caption={src/styles/Header.css}]
.header {
  padding: 2rem 0;
}

.header__name {
  font-size: 2rem;
  font-weight: 400;
  color: #8B4513;
  margin-bottom: 0.5rem;
}

.header__tagline {
  font-size: 0.9rem;
  color: #666;
  margin-bottom: 1rem;
}

.header__obsessions {
  font-size: 0.9rem;
  margin-bottom: 0.5rem;
}

.header__belief {
  font-size: 0.9rem;
}

.header__link {
  color: inherit;
  text-decoration: underline;
}

.header__link:hover {
  text-decoration: none;
}
\end{lstlisting}

\begin{explanation}
\textbf{Line 2:} \texttt{padding: 2rem 0} - 2rem top/bottom, 0 left/right. \texttt{rem} is a unit relative to root font size (usually 16px, so 2rem = 32px).

\textbf{Line 8:} \texttt{\#8B4513} - Hex color code for a brownish shade (matching your design's name color).

\textbf{Line 28:} \texttt{color: inherit} - Use the same color as parent element.

\textbf{Line 32:} \texttt{:hover} - Pseudo-class that applies when mouse hovers over element.
\end{explanation}

\subsection{Step 10.3: Import Header into App}

Update \texttt{src/App.tsx}:

\begin{lstlisting}[language=JavaScript,caption={src/App.tsx with Header}]
import Header from './components/Header';
import './styles/Header.css';

function App() {
  return (
    <main className="container">
      <Header />
    </main>
  );
}

export default App;
\end{lstlisting}

\begin{explanation}
\textbf{Line 1:} Import the Header component.

\textbf{Line 2:} Import its styles.

\textbf{Line 7:} \texttt{<Header />} - Use the component. Self-closing tag since it has no children.
\end{explanation}

Add container styles to \texttt{src/index.css}:

\begin{lstlisting}[language=CSS,caption={Add to src/index.css}]
.container {
  max-width: 800px;
  margin: 0 auto;
  padding: 2rem;
}
\end{lstlisting}

\begin{explanation}
\textbf{Line 2:} \texttt{max-width: 800px} - Content won't exceed 800px wide.

\textbf{Line 3:} \texttt{margin: 0 auto} - 0 top/bottom margin, \texttt{auto} left/right centers the element.

\textbf{Line 4:} \texttt{padding: 2rem} - Space inside the container.
\end{explanation}

\begin{checkpoint}
Your browser should now show:
\begin{itemize}
    \item[$\square$] Your name in brownish color
    \item[$\square$] Tagline below it
    \item[$\square$] Obsessions line
    \item[$\square$] Belief statement with underlined link
    \item[$\square$] Content centered on the page
\end{itemize}
\end{checkpoint}

\begin{gitcommit}
\begin{lstlisting}[language=bash]
git add .
git commit -m "Add Header component with styles"
git push
\end{lstlisting}
\end{gitcommit}

%==============================================================================
% CHAPTER 11: EXPERIENCE SECTION
%==============================================================================
\section{Chapter 11: Building the Experience Section}

\subsection{Step 11.1: Create Section Component}

A reusable Section component will wrap each content section.

Create \texttt{src/components/Section.tsx}:

\begin{lstlisting}[language=JavaScript,caption={src/components/Section.tsx}]
import { ReactNode } from 'react';

interface SectionProps {
  title: string;
  children: ReactNode;
}

function Section({ title, children }: SectionProps) {
  return (
    <section className="section">
      <h2 className="section__title">{title}</h2>
      <div className="section__content">
        {children}
      </div>
    </section>
  );
}

export default Section;
\end{lstlisting}

\begin{explanation}
\textbf{Line 1:} \texttt{ReactNode} is a TypeScript type that means ``anything React can render''---text, elements, components, arrays, null, etc.

\textbf{Lines 3-6:} Props interface:
\begin{itemize}
    \item \texttt{title: string} - The section heading
    \item \texttt{children: ReactNode} - Whatever is placed between \texttt{<Section>} and \texttt{</Section>}
\end{itemize}

\textbf{Line 8:} Destructure both props directly in the function signature.

\textbf{Line 13:} \texttt{\{children\}} - Render whatever was passed as children.
\end{explanation}

\begin{concept}
\textbf{The Children Prop}

In React, components can wrap other content:
\begin{lstlisting}[language=JavaScript]
<Section title="work">
  <ExperienceCard ... />
  <ExperienceCard ... />
</Section>
\end{lstlisting}

Everything between the opening and closing tags becomes \texttt{children}. This is how you create reusable wrapper components.
\end{concept}

\subsection{Step 11.2: Create ExperienceCard Component}

Create \texttt{src/components/ExperienceCard.tsx}:

\begin{lstlisting}[language=JavaScript,caption={src/components/ExperienceCard.tsx}]
import { Experience } from '../data/content';

interface ExperienceCardProps {
  experience: Experience;
}

function ExperienceCard({ experience }: ExperienceCardProps) {
  const { title, company, period, description, tags } = experience;

  return (
    <article className="experience-card">
      <div className="experience-card__header">
        <h3 className="experience-card__title">
          {title}, {company}
        </h3>
        {tags && (
          <div className="experience-card__tags">
            {tags.map((tag, index) => (
              <span key={index} className="experience-card__tag">
                {tag}
              </span>
            ))}
          </div>
        )}
      </div>
      <p className="experience-card__period">{period}</p>
      <p className="experience-card__description">{description}</p>
    </article>
  );
}

export default ExperienceCard;
\end{lstlisting}

\begin{explanation}
\textbf{Line 1:} Import the Experience interface (for type checking).

\textbf{Lines 3-5:} This component expects an \texttt{experience} prop of type \texttt{Experience}.

\textbf{Line 11:} \texttt{<article>} - Semantic HTML for self-contained content.

\textbf{Line 16:} \texttt{tags \&\& (...)} - Conditional rendering. If \texttt{tags} exists (is truthy), render what's in parentheses. If \texttt{tags} is undefined, nothing renders.

\textbf{Lines 18-22:} \texttt{tags.map((tag, index) => ...)} - The \texttt{map} function transforms each array element:
\begin{itemize}
    \item \texttt{tag} - Current tag string
    \item \texttt{index} - Position in array (0, 1, 2...)
    \item Returns JSX for each tag
\end{itemize}

\textbf{Line 19:} \texttt{key=\{index\}} - React requires a unique \texttt{key} for each item in a list. This helps React efficiently update the DOM.
\end{explanation}

\begin{concept}
\textbf{Array.map() - Transforming Arrays}

\texttt{map} is one of the most important functions in React:

\begin{lstlisting}[language=JavaScript]
const numbers = [1, 2, 3];
const doubled = numbers.map(n => n * 2);
// doubled = [2, 4, 6]

const names = ['Alice', 'Bob'];
const greetings = names.map(name => `Hello, ${name}!`);
// greetings = ['Hello, Alice!', 'Hello, Bob!']
\end{lstlisting}

In React, we use \texttt{map} to convert data arrays into arrays of components.
\end{concept}

\subsection{Step 11.3: Add Styles}

Create \texttt{src/styles/Section.css}:

\begin{lstlisting}[language=CSS,caption={src/styles/Section.css}]
.section {
  margin: 3rem 0;
}

.section__title {
  font-size: 1.25rem;
  font-weight: 400;
  color: #333;
  margin-bottom: 1.5rem;
  padding-bottom: 0.5rem;
  border-bottom: 1px solid #e0e0e0;
}

.section__content {
  display: flex;
  flex-direction: column;
  gap: 2rem;
}
\end{lstlisting}

\begin{explanation}
\textbf{Line 15:} \texttt{display: flex} - Enable Flexbox layout.

\textbf{Line 16:} \texttt{flex-direction: column} - Stack children vertically.

\textbf{Line 17:} \texttt{gap: 2rem} - Space between children (modern CSS, much cleaner than margins).
\end{explanation}

\begin{concept}
\textbf{Flexbox}

Flexbox is a CSS layout system. Key concepts:
\begin{itemize}
    \item \texttt{display: flex} - Enable on container
    \item \texttt{flex-direction} - row (horizontal) or column (vertical)
    \item \texttt{justify-content} - Alignment along main axis
    \item \texttt{align-items} - Alignment along cross axis
    \item \texttt{gap} - Space between items
\end{itemize}
\end{concept}

Create \texttt{src/styles/ExperienceCard.css}:

\begin{lstlisting}[language=CSS,caption={src/styles/ExperienceCard.css}]
.experience-card {
  /* Card is just a logical grouping, no visible styling */
}

.experience-card__header {
  display: flex;
  align-items: center;
  gap: 0.75rem;
  flex-wrap: wrap;
}

.experience-card__title {
  font-size: 1rem;
  font-weight: 500;
  color: #333;
  margin: 0;
}

.experience-card__tags {
  display: flex;
  gap: 0.5rem;
}

.experience-card__tag {
  font-size: 0.7rem;
  padding: 0.15rem 0.4rem;
  background: #333;
  color: white;
  border-radius: 3px;
}

.experience-card__period {
  font-size: 0.8rem;
  color: #888;
  font-style: italic;
  margin: 0.25rem 0 0.5rem 0;
}

.experience-card__description {
  font-size: 0.9rem;
  color: #555;
  line-height: 1.5;
  margin: 0;
}
\end{lstlisting}

\subsection{Step 11.4: Integrate into App}

Update \texttt{src/App.tsx}:

\begin{lstlisting}[language=JavaScript,caption={src/App.tsx with Experience section}]
import Header from './components/Header';
import Section from './components/Section';
import ExperienceCard from './components/ExperienceCard';
import { experiences } from './data/content';

import './styles/Header.css';
import './styles/Section.css';
import './styles/ExperienceCard.css';

function App() {
  return (
    <main className="container">
      <Header />

      <Section title="work">
        {experiences.map((exp, index) => (
          <ExperienceCard key={index} experience={exp} />
        ))}
      </Section>
    </main>
  );
}

export default App;
\end{lstlisting}

\begin{explanation}
\textbf{Lines 2-4:} Import all our new components and data.

\textbf{Lines 6-8:} Import all CSS files.

\textbf{Lines 15-19:} Use the Section component with title ``work,'' then map over experiences to create ExperienceCards.
\end{explanation}

\begin{checkpoint}
Your page should now show:
\begin{itemize}
    \item[$\square$] Header section (name, tagline, obsessions, belief)
    \item[$\square$] ``work'' section heading with underline
    \item[$\square$] Three experience cards with title, tags, period, and description
\end{itemize}
\end{checkpoint}

\begin{gitcommit}
\begin{lstlisting}[language=bash]
git add .
git commit -m "Add Section and ExperienceCard components"
git push
\end{lstlisting}
\end{gitcommit}

%==============================================================================
% CHAPTER 12: RESEARCH AND PROJECTS SECTIONS
%==============================================================================
\section{Chapter 12: Research and Projects Sections}

Following the same pattern, let's add the remaining sections.

\subsection{Step 12.1: Create ResearchCard Component}

Create \texttt{src/components/ResearchCard.tsx}:

\begin{lstlisting}[language=JavaScript,caption={src/components/ResearchCard.tsx}]
import { Research } from '../data/content';

interface ResearchCardProps {
  research: Research;
}

function ResearchCard({ research }: ResearchCardProps) {
  const { title, venue, note, description } = research;

  return (
    <article className="research-card">
      <h3 className="research-card__title">{title}</h3>
      {venue && (
        <p className="research-card__venue">
          <em>{note ? `${note} at ${venue}` : venue}</em>
        </p>
      )}
      {description && (
        <p className="research-card__description">{description}</p>
      )}
    </article>
  );
}

export default ResearchCard;
\end{lstlisting}

\begin{explanation}
\textbf{Line 15:} Template literal with conditional: If \texttt{note} exists, show ``note at venue'', otherwise just venue.

\textbf{Lines 13-20:} Multiple conditional renders---only show elements if data exists.
\end{explanation}

\subsection{Step 12.2: Create ProjectCard Component}

Create \texttt{src/components/ProjectCard.tsx}:

\begin{lstlisting}[language=JavaScript,caption={src/components/ProjectCard.tsx}]
import { Project } from '../data/content';

interface ProjectCardProps {
  project: Project;
}

function ProjectCard({ project }: ProjectCardProps) {
  const { title, tagline, metrics } = project;

  return (
    <article className="project-card">
      <h3 className="project-card__title">Founder, {title}</h3>
      <p className="project-card__tagline"><em>{tagline}</em></p>
      {metrics && (
        <p className="project-card__metrics">{metrics}</p>
      )}
    </article>
  );
}

export default ProjectCard;
\end{lstlisting}

\subsection{Step 12.3: Add Styles}

Create \texttt{src/styles/ResearchCard.css}:

\begin{lstlisting}[language=CSS,caption={src/styles/ResearchCard.css}]
.research-card__title {
  font-size: 1rem;
  font-weight: 500;
  color: #333;
  margin: 0 0 0.25rem 0;
}

.research-card__venue {
  font-size: 0.85rem;
  color: #666;
  margin: 0 0 0.5rem 0;
}

.research-card__description {
  font-size: 0.9rem;
  color: #555;
  line-height: 1.5;
  margin: 0;
}
\end{lstlisting}

Create \texttt{src/styles/ProjectCard.css}:

\begin{lstlisting}[language=CSS,caption={src/styles/ProjectCard.css}]
.project-card__title {
  font-size: 1rem;
  font-weight: 500;
  color: #333;
  margin: 0;
}

.project-card__tagline {
  font-size: 0.9rem;
  color: #555;
  margin: 0.25rem 0;
}

.project-card__metrics {
  font-size: 0.85rem;
  color: #666;
  margin: 0;
}
\end{lstlisting}

\subsection{Step 12.4: Update App.tsx}

\begin{lstlisting}[language=JavaScript,caption={src/App.tsx - Complete version}]
import Header from './components/Header';
import Section from './components/Section';
import ExperienceCard from './components/ExperienceCard';
import ResearchCard from './components/ResearchCard';
import ProjectCard from './components/ProjectCard';
import { experiences, research, projects } from './data/content';

import './styles/Header.css';
import './styles/Section.css';
import './styles/ExperienceCard.css';
import './styles/ResearchCard.css';
import './styles/ProjectCard.css';

function App() {
  return (
    <main className="container">
      <Header />

      <Section title="work">
        {experiences.map((exp, index) => (
          <ExperienceCard key={index} experience={exp} />
        ))}
      </Section>

      <Section title="research">
        {research.map((item, index) => (
          <ResearchCard key={index} research={item} />
        ))}
      </Section>

      <Section title="other">
        {projects.map((proj, index) => (
          <ProjectCard key={index} project={proj} />
        ))}
      </Section>
    </main>
  );
}

export default App;
\end{lstlisting}

\begin{checkpoint}
Your page should now display:
\begin{itemize}
    \item[$\square$] Header
    \item[$\square$] Work section with 3 experience cards
    \item[$\square$] Research section with 2 research cards
    \item[$\square$] Other section with 2 project cards
\end{itemize}
\end{checkpoint}

\begin{gitcommit}
\begin{lstlisting}[language=bash]
git add .
git commit -m "Add Research and Project sections with cards"
git push
\end{lstlisting}
\end{gitcommit}

%==============================================================================
% CHAPTER 13: SOCIAL LINKS FOOTER
%==============================================================================
\section{Chapter 13: Social Links Footer}

\subsection{Step 13.1: Create SVG Icons}

For this section, you'll export SVG icons from Figma. For now, let's create placeholder icons using simple shapes.

Create \texttt{src/components/SocialLinks.tsx}:

\begin{lstlisting}[language=JavaScript,caption={src/components/SocialLinks.tsx}]
import { socialLinks } from '../data/content';

function SocialLinks() {
  return (
    <footer className="social-links">
      <a
        href={socialLinks.github}
        target="_blank"
        rel="noopener noreferrer"
        className="social-links__link"
        aria-label="GitHub"
      >
        <svg width="20" height="20" viewBox="0 0 24 24" fill="currentColor">
          <path d="M12 0c-6.626 0-12 5.373-12 12 0 5.302 3.438 9.8 8.207 11.387.599.111.793-.261.793-.577v-2.234c-3.338.726-4.033-1.416-4.033-1.416-.546-1.387-1.333-1.756-1.333-1.756-1.089-.745.083-.729.083-.729 1.205.084 1.839 1.237 1.839 1.237 1.07 1.834 2.807 1.304 3.492.997.107-.775.418-1.305.762-1.604-2.665-.305-5.467-1.334-5.467-5.931 0-1.311.469-2.381 1.236-3.221-.124-.303-.535-1.524.117-3.176 0 0 1.008-.322 3.301 1.23.957-.266 1.983-.399 3.003-.404 1.02.005 2.047.138 3.006.404 2.291-1.552 3.297-1.23 3.297-1.23.653 1.653.242 2.874.118 3.176.77.84 1.235 1.911 1.235 3.221 0 4.609-2.807 5.624-5.479 5.921.43.372.823 1.102.823 2.222v3.293c0 .319.192.694.801.576 4.765-1.589 8.199-6.086 8.199-11.386 0-6.627-5.373-12-12-12z"/>
        </svg>
      </a>

      <a
        href={socialLinks.twitter}
        target="_blank"
        rel="noopener noreferrer"
        className="social-links__link"
        aria-label="Twitter"
      >
        <svg width="20" height="20" viewBox="0 0 24 24" fill="currentColor">
          <path d="M18.244 2.25h3.308l-7.227 8.26 8.502 11.24H16.17l-5.214-6.817L4.99 21.75H1.68l7.73-8.835L1.254 2.25H8.08l4.713 6.231zm-1.161 17.52h1.833L7.084 4.126H5.117z"/>
        </svg>
      </a>

      <a
        href={socialLinks.linkedin}
        target="_blank"
        rel="noopener noreferrer"
        className="social-links__link"
        aria-label="LinkedIn"
      >
        <svg width="20" height="20" viewBox="0 0 24 24" fill="currentColor">
          <path d="M20.447 20.452h-3.554v-5.569c0-1.328-.027-3.037-1.852-3.037-1.853 0-2.136 1.445-2.136 2.939v5.667H9.351V9h3.414v1.561h.046c.477-.9 1.637-1.85 3.37-1.85 3.601 0 4.267 2.37 4.267 5.455v6.286zM5.337 7.433c-1.144 0-2.063-.926-2.063-2.065 0-1.138.92-2.063 2.063-2.063 1.14 0 2.064.925 2.064 2.063 0 1.139-.925 2.065-2.064 2.065zm1.782 13.019H3.555V9h3.564v11.452zM22.225 0H1.771C.792 0 0 .774 0 1.729v20.542C0 23.227.792 24 1.771 24h20.451C23.2 24 24 23.227 24 22.271V1.729C24 .774 23.2 0 22.222 0h.003z"/>
        </svg>
      </a>

      <a
        href={socialLinks.email}
        className="social-links__link"
        aria-label="Email"
      >
        <svg width="20" height="20" viewBox="0 0 24 24" fill="currentColor">
          <path d="M20 4H4c-1.1 0-1.99.9-1.99 2L2 18c0 1.1.9 2 2 2h16c1.1 0 2-.9 2-2V6c0-1.1-.9-2-2-2zm0 4l-8 5-8-5V6l8 5 8-5v2z"/>
        </svg>
      </a>
    </footer>
  );
}

export default SocialLinks;
\end{lstlisting}

\begin{explanation}
\textbf{Line 8-9:} \texttt{target="\_blank"} opens link in new tab. \texttt{rel="noopener noreferrer"} is a security measure that prevents the new page from accessing your page.

\textbf{Line 11:} \texttt{aria-label} provides text for screen readers since the link only contains an icon.

\textbf{Lines 13-15:} Inline SVG for the GitHub icon. \texttt{fill="currentColor"} makes the icon inherit its parent's text color.
\end{explanation}

\begin{concept}
\textbf{SVG (Scalable Vector Graphics)}

SVG is a format for graphics defined by math (paths, shapes) rather than pixels. Benefits:
\begin{itemize}
    \item Scales to any size without losing quality
    \item Small file size for simple shapes
    \item Can be styled with CSS
    \item Can be embedded directly in HTML/JSX
\end{itemize}

You'll export SVGs from Figma for your actual icons.
\end{concept}

Create \texttt{src/styles/SocialLinks.css}:

\begin{lstlisting}[language=CSS,caption={src/styles/SocialLinks.css}]
.social-links {
  display: flex;
  gap: 1rem;
  padding: 2rem 0;
  margin-top: 2rem;
  border-top: 1px solid #e0e0e0;
}

.social-links__link {
  color: #666;
  transition: color 0.2s ease;
}

.social-links__link:hover {
  color: #333;
}
\end{lstlisting}

\begin{explanation}
\textbf{Line 11:} \texttt{transition: color 0.2s ease} - Animate the color change over 0.2 seconds with an ``ease'' timing function (starts slow, speeds up, slows down).
\end{explanation}

Update \texttt{src/App.tsx} to include SocialLinks:

\begin{lstlisting}[language=JavaScript,caption={Add to App.tsx}]
import SocialLinks from './components/SocialLinks';
import './styles/SocialLinks.css';

// ... in the return statement, after the last Section:
      <SocialLinks />
    </main>
\end{lstlisting}

\begin{checkpoint}
Your page should now show:
\begin{itemize}
    \item[$\square$] All previous sections
    \item[$\square$] Footer with 4 social icons
    \item[$\square$] Icons change color on hover
\end{itemize}
\end{checkpoint}

\begin{gitcommit}
\begin{lstlisting}[language=bash]
git add .
git commit -m "Add SocialLinks footer with SVG icons"
git push
\end{lstlisting}
\end{gitcommit}

%==============================================================================
% CHAPTER 14: DARK MODE
%==============================================================================
\section{Chapter 14: Implementing Dark Mode}

Dark mode requires understanding \textbf{state} and \textbf{CSS custom properties}.

\subsection{Step 14.1: Create Theme Context}

React \textbf{Context} lets us share data across components without passing props through every level.

Create \texttt{src/context/ThemeContext.tsx}:

\begin{lstlisting}[language=JavaScript,caption={src/context/ThemeContext.tsx}]
import { createContext, useContext, useState, useEffect, ReactNode } from 'react';

type Theme = 'light' | 'dark';

interface ThemeContextType {
  theme: Theme;
  toggleTheme: () => void;
}

const ThemeContext = createContext<ThemeContextType | undefined>(undefined);

export function ThemeProvider({ children }: { children: ReactNode }) {
  const [theme, setTheme] = useState<Theme>(() => {
    // Check localStorage for saved preference
    const saved = localStorage.getItem('theme') as Theme;
    if (saved) return saved;

    // Check system preference
    if (window.matchMedia('(prefers-color-scheme: dark)').matches) {
      return 'dark';
    }
    return 'light';
  });

  useEffect(() => {
    // Update document attribute and localStorage when theme changes
    document.documentElement.setAttribute('data-theme', theme);
    localStorage.setItem('theme', theme);
  }, [theme]);

  const toggleTheme = () => {
    setTheme(prev => prev === 'light' ? 'dark' : 'light');
  };

  return (
    <ThemeContext.Provider value={{ theme, toggleTheme }}>
      {children}
    </ThemeContext.Provider>
  );
}

export function useTheme() {
  const context = useContext(ThemeContext);
  if (context === undefined) {
    throw new Error('useTheme must be used within a ThemeProvider');
  }
  return context;
}
\end{lstlisting}

\begin{explanation}
\textbf{Line 3:} \texttt{type Theme = 'light' | 'dark'} - A ``union type.'' Theme can ONLY be 'light' or 'dark', nothing else.

\textbf{Line 10:} \texttt{createContext} creates a context object. The type parameter \texttt{ThemeContextType | undefined} means it can be our type OR undefined (before a provider wraps it).

\textbf{Lines 13-23:} Initialize state with a function that:
\begin{enumerate}
    \item Checks localStorage for a saved preference
    \item Falls back to system preference (\texttt{prefers-color-scheme})
    \item Defaults to light
\end{enumerate}

\textbf{Lines 25-28:} \texttt{useEffect} runs side effects. This runs whenever \texttt{theme} changes (that's what \texttt{[theme]} means---the ``dependency array'').

\textbf{Line 26:} Set a data attribute on the HTML element. We'll use this in CSS.

\textbf{Lines 30-32:} Toggle function---if light, become dark; if dark, become light.

\textbf{Lines 34-38:} \texttt{Provider} makes the context available to all children.

\textbf{Lines 41-46:} Custom hook for easy access. Throws error if used outside provider (helps catch bugs).
\end{explanation}

\begin{concept}
\textbf{useEffect Hook}

\texttt{useEffect} handles ``side effects''---things that affect the outside world:
\begin{itemize}
    \item DOM manipulation
    \item API calls
    \item localStorage
    \item Event listeners
\end{itemize}

Syntax:
\begin{lstlisting}[language=JavaScript]
useEffect(() => {
  // Effect code
}, [dependency1, dependency2]);
\end{lstlisting}

The effect runs:
\begin{enumerate}
    \item After initial render
    \item After any render where dependencies changed
\end{enumerate}
\end{concept}

\subsection{Step 14.2: Create ThemeToggle Component}

Create \texttt{src/components/ThemeToggle.tsx}:

\begin{lstlisting}[language=JavaScript,caption={src/components/ThemeToggle.tsx}]
import { useTheme } from '../context/ThemeContext';

function ThemeToggle() {
  const { theme, toggleTheme } = useTheme();

  return (
    <button
      className="theme-toggle"
      onClick={toggleTheme}
      aria-label={`Switch to ${theme === 'light' ? 'dark' : 'light'} mode`}
    >
      {theme === 'light' ? (
        <svg width="20" height="20" viewBox="0 0 24 24" fill="currentColor">
          <path d="M21 12.79A9 9 0 1 1 11.21 3 7 7 0 0 0 21 12.79z"/>
        </svg>
      ) : (
        <svg width="20" height="20" viewBox="0 0 24 24" fill="currentColor">
          <circle cx="12" cy="12" r="5"/>
          <line x1="12" y1="1" x2="12" y2="3" stroke="currentColor" strokeWidth="2"/>
          <line x1="12" y1="21" x2="12" y2="23" stroke="currentColor" strokeWidth="2"/>
          <line x1="4.22" y1="4.22" x2="5.64" y2="5.64" stroke="currentColor" strokeWidth="2"/>
          <line x1="18.36" y1="18.36" x2="19.78" y2="19.78" stroke="currentColor" strokeWidth="2"/>
          <line x1="1" y1="12" x2="3" y2="12" stroke="currentColor" strokeWidth="2"/>
          <line x1="21" y1="12" x2="23" y2="12" stroke="currentColor" strokeWidth="2"/>
          <line x1="4.22" y1="19.78" x2="5.64" y2="18.36" stroke="currentColor" strokeWidth="2"/>
          <line x1="18.36" y1="5.64" x2="19.78" y2="4.22" stroke="currentColor" strokeWidth="2"/>
        </svg>
      )}
    </button>
  );
}

export default ThemeToggle;
\end{lstlisting}

\begin{explanation}
\textbf{Line 4:} Use our custom hook to get current theme and toggle function.

\textbf{Lines 12-27:} Conditional rendering---show moon icon in light mode, sun icon in dark mode.
\end{explanation}

\subsection{Step 14.3: Update CSS for Theming}

Update \texttt{src/index.css}:

\begin{lstlisting}[language=CSS,caption={src/index.css with CSS custom properties}]
:root {
  --color-bg: #fafafa;
  --color-text: #333;
  --color-text-secondary: #666;
  --color-text-muted: #888;
  --color-accent: #8B4513;
  --color-border: #e0e0e0;
  --color-tag-bg: #333;
  --color-tag-text: white;
}

[data-theme='dark'] {
  --color-bg: #1a1a1a;
  --color-text: #e0e0e0;
  --color-text-secondary: #a0a0a0;
  --color-text-muted: #808080;
  --color-accent: #d4956a;
  --color-border: #333;
  --color-tag-bg: #e0e0e0;
  --color-tag-text: #1a1a1a;
}

* {
  margin: 0;
  padding: 0;
  box-sizing: border-box;
}

body {
  font-family: -apple-system, BlinkMacSystemFont, 'Segoe UI', Roboto,
    Oxygen, Ubuntu, Cantarell, sans-serif;
  line-height: 1.6;
  color: var(--color-text);
  background-color: var(--color-bg);
  transition: background-color 0.3s ease, color 0.3s ease;
}

.container {
  max-width: 800px;
  margin: 0 auto;
  padding: 2rem;
}
\end{lstlisting}

\begin{explanation}
\textbf{Lines 1-10:} \texttt{:root} is the HTML element. We define CSS custom properties (variables) with \texttt{--name} syntax.

\textbf{Lines 12-21:} When \texttt{data-theme='dark'} is set on HTML element, override the variables with dark theme colors.

\textbf{Line 33:} \texttt{var(--color-text)} uses the variable. In light mode it's \#333, in dark mode it's \#e0e0e0.

\textbf{Line 35:} Smooth transition when theme changes.
\end{explanation}

\begin{concept}
\textbf{CSS Custom Properties (Variables)}

Variables in CSS:
\begin{lstlisting}[language=CSS]
:root {
  --my-color: blue;
}

.element {
  color: var(--my-color);  /* Uses blue */
}
\end{lstlisting}

Benefits:
\begin{itemize}
    \item Change one value, affects everywhere it's used
    \item Can be overridden in specific contexts (like dark mode)
    \item Computed at runtime (unlike preprocessor variables)
\end{itemize}
\end{concept}

Now update all component CSS files to use these variables. For example, in \texttt{Header.css}:

\begin{lstlisting}[language=CSS,caption={Update Header.css}]
.header__name {
  font-size: 2rem;
  font-weight: 400;
  color: var(--color-accent);
  margin-bottom: 0.5rem;
}

.header__tagline {
  font-size: 0.9rem;
  color: var(--color-text-secondary);
  margin-bottom: 1rem;
}
/* ... update all color values to use variables */
\end{lstlisting}

\subsection{Step 14.4: Wrap App with Provider}

Create folder and file \texttt{src/context/ThemeContext.tsx} (we did this above).

Update \texttt{src/main.tsx}:

\begin{lstlisting}[language=JavaScript,caption={src/main.tsx with ThemeProvider}]
import React from 'react'
import ReactDOM from 'react-dom/client'
import App from './App.tsx'
import { ThemeProvider } from './context/ThemeContext'
import './index.css'

ReactDOM.createRoot(document.getElementById('root')!).render(
  <React.StrictMode>
    <ThemeProvider>
      <App />
    </ThemeProvider>
  </React.StrictMode>,
)
\end{lstlisting}

Add ThemeToggle to App.tsx header area:

\begin{lstlisting}[language=JavaScript,caption={Add ThemeToggle to layout}]
import ThemeToggle from './components/ThemeToggle';
import './styles/ThemeToggle.css';

// In the return, at the top:
<div className="top-bar">
  <span>other</span>
  <ThemeToggle />
</div>
\end{lstlisting}

\begin{checkpoint}
\begin{itemize}
    \item[$\square$] Click the theme toggle button
    \item[$\square$] Page should smoothly transition between light and dark
    \item[$\square$] Refresh the page---your preference should be remembered
    \item[$\square$] All text should be readable in both modes
\end{itemize}
\end{checkpoint}

\begin{gitcommit}
\begin{lstlisting}[language=bash]
git add .
git commit -m "Implement dark mode with ThemeContext and CSS variables"
git push
\end{lstlisting}
\end{gitcommit}

%==============================================================================
% CHAPTER 15: SKILLS GRAPH VISUALIZATION
%==============================================================================
\section{Chapter 15: Skills Graph Visualization}

The skills graph is the most complex component. We'll create an interactive visualization.

\begin{concept}
\textbf{Canvas vs SVG}

For graphics in web:
\begin{itemize}
    \item \textbf{SVG}: Vector-based, each element is a DOM node, good for <1000 elements
    \item \textbf{Canvas}: Pixel-based, one DOM element, good for animations and many elements
\end{itemize}

For your skills graph, SVG is appropriate since we have few elements.
\end{concept}

\subsection{Step 15.1: Create SkillsGraph Component}

Create \texttt{src/components/SkillsGraph.tsx}:

\begin{lstlisting}[language=JavaScript,caption={src/components/SkillsGraph.tsx}]
import { useEffect, useRef, useState } from 'react';
import { skills } from '../data/content';

interface Point {
  x: number;
  y: number;
  vx: number;
  vy: number;
}

function SkillsGraph() {
  const containerRef = useRef<HTMLDivElement>(null);
  const [points, setPoints] = useState<Point[]>([]);

  useEffect(() => {
    // Generate random points for the graph
    const generatePoints = () => {
      const newPoints: Point[] = [];
      for (let i = 0; i < 50; i++) {
        newPoints.push({
          x: Math.random() * 200,
          y: Math.random() * 150,
          vx: (Math.random() - 0.5) * 0.5,
          vy: (Math.random() - 0.5) * 0.5
        });
      }
      return newPoints;
    };

    setPoints(generatePoints());
  }, []);

  // Animation loop
  useEffect(() => {
    if (points.length === 0) return;

    const animate = () => {
      setPoints(prevPoints =>
        prevPoints.map(point => {
          let newX = point.x + point.vx;
          let newY = point.y + point.vy;
          let newVx = point.vx;
          let newVy = point.vy;

          // Bounce off walls
          if (newX < 0 || newX > 200) newVx *= -1;
          if (newY < 0 || newY > 150) newVy *= -1;

          return {
            x: Math.max(0, Math.min(200, newX)),
            y: Math.max(0, Math.min(150, newY)),
            vx: newVx,
            vy: newVy
          };
        })
      );
    };

    const intervalId = setInterval(animate, 50);
    return () => clearInterval(intervalId);
  }, [points.length]);

  return (
    <div className="skills-graph" ref={containerRef}>
      <div className="skills-graph__header">
        Graph View of current skills
      </div>
      <svg
        className="skills-graph__svg"
        viewBox="0 0 200 150"
        preserveAspectRatio="xMidYMid meet"
      >
        {/* Draw connecting lines between nearby points */}
        {points.map((point, i) =>
          points.slice(i + 1).map((other, j) => {
            const distance = Math.sqrt(
              Math.pow(point.x - other.x, 2) +
              Math.pow(point.y - other.y, 2)
            );
            if (distance < 40) {
              return (
                <line
                  key={`line-${i}-${j}`}
                  x1={point.x}
                  y1={point.y}
                  x2={other.x}
                  y2={other.y}
                  stroke="var(--color-text-muted)"
                  strokeWidth="0.5"
                  opacity={1 - distance / 40}
                />
              );
            }
            return null;
          })
        )}

        {/* Draw points */}
        {points.map((point, i) => (
          <circle
            key={`point-${i}`}
            cx={point.x}
            cy={point.y}
            r="2"
            fill="var(--color-text)"
          />
        ))}
      </svg>

      <ul className="skills-graph__list">
        {skills.map((skill, index) => (
          <li key={index} className="skills-graph__skill">
            {skill}
          </li>
        ))}
      </ul>
    </div>
  );
}

export default SkillsGraph;
\end{lstlisting}

\begin{explanation}
\textbf{Line 12:} \texttt{useRef} creates a reference to a DOM element. \texttt{<HTMLDivElement>} is the TypeScript type.

\textbf{Lines 15-30:} First useEffect generates initial random points when component mounts.

\textbf{Lines 33-59:} Second useEffect runs animation:
\begin{itemize}
    \item Updates point positions based on velocity
    \item Reverses velocity at boundaries (bounce effect)
    \item \texttt{setInterval} calls \texttt{animate} every 50ms
    \item Return function cleans up interval when component unmounts
\end{itemize}

\textbf{Lines 71-89:} Draw lines between nearby points. The nested map compares each point with every other point.

\textbf{Line 87:} \texttt{opacity=\{1 - distance / 40\}} - Lines fade as distance increases.
\end{explanation}

\begin{concept}
\textbf{useRef Hook}

\texttt{useRef} has two main uses:
\begin{enumerate}
    \item Reference DOM elements directly
    \item Store mutable values that don't trigger re-renders
\end{enumerate}

\begin{lstlisting}[language=JavaScript]
const myRef = useRef<HTMLDivElement>(null);
// Access the DOM element: myRef.current
\end{lstlisting}
\end{concept}

\begin{gitcommit}
\begin{lstlisting}[language=bash]
git add .
git commit -m "Add animated SkillsGraph visualization component"
git push
\end{lstlisting}
\end{gitcommit}

%==============================================================================
% CHAPTER 16: RESPONSIVE DESIGN
%==============================================================================
\section{Chapter 16: Responsive Design}

Your website must look good on all screen sizes.

\begin{concept}
\textbf{Media Queries}

CSS media queries apply styles conditionally based on screen size:
\begin{lstlisting}[language=CSS]
/* Default styles (mobile first) */
.container {
  padding: 1rem;
}

/* Larger screens */
@media (min-width: 768px) {
  .container {
    padding: 2rem;
  }
}
\end{lstlisting}

Common breakpoints:
\begin{itemize}
    \item 480px - Small phones
    \item 768px - Tablets
    \item 1024px - Laptops
    \item 1200px - Desktops
\end{itemize}
\end{concept}

Add responsive styles to your CSS files. Example additions:

\begin{lstlisting}[language=CSS,caption={Responsive additions}]
/* Mobile-first base styles */
.header__name {
  font-size: 1.5rem;
}

@media (min-width: 768px) {
  .header__name {
    font-size: 2rem;
  }
}

/* Layout adjustments */
@media (max-width: 768px) {
  .top-bar {
    flex-direction: column;
    align-items: flex-start;
  }

  .skills-graph {
    margin-top: 2rem;
  }
}
\end{lstlisting}

\begin{checkpoint}
Test responsiveness:
\begin{itemize}
    \item[$\square$] Open browser DevTools (F12 or right-click → Inspect)
    \item[$\square$] Click the device toggle icon (or Ctrl+Shift+M)
    \item[$\square$] Test at different screen sizes
    \item[$\square$] Ensure text is readable and layout makes sense at all sizes
\end{itemize}
\end{checkpoint}

\begin{gitcommit}
\begin{lstlisting}[language=bash]
git add .
git commit -m "Add responsive design with media queries"
git push
\end{lstlisting}
\end{gitcommit}

%==============================================================================
% CHAPTER 17: BUILDING FOR PRODUCTION
%==============================================================================
\section{Chapter 17: Building for Production}

Your development server is for development. For the real internet, you need a production build.

\subsection{Step 17.1: Create Production Build}

\begin{lstlisting}[language=bash,caption={Build command}]
npm run build
\end{lstlisting}

\begin{explanation}
This command:
\begin{enumerate}
    \item Compiles TypeScript to JavaScript
    \item Bundles all files together
    \item Minifies code (removes whitespace, shortens variable names)
    \item Outputs to the \texttt{dist} folder
\end{enumerate}
\end{explanation}

\subsection{Step 17.2: Preview Production Build}

\begin{lstlisting}[language=bash,caption={Preview production build}]
npm run preview
\end{lstlisting}

This starts a local server serving your production build.

\begin{checkpoint}
\begin{itemize}
    \item[$\square$] Visit the preview URL (usually http://localhost:4173)
    \item[$\square$] Verify everything works exactly like development
    \item[$\square$] Check that dark mode persists
    \item[$\square$] Check that animations work
\end{itemize}
\end{checkpoint}

\subsection{Step 17.3: Deploy to GitHub Pages (Free Hosting)}

Update \texttt{vite.config.ts}:

\begin{lstlisting}[language=JavaScript,caption={vite.config.ts for GitHub Pages}]
import { defineConfig } from 'vite'
import react from '@vitejs/plugin-react'

export default defineConfig({
  plugins: [react()],
  base: '/my-portfolio/', // Replace with your repo name
})
\end{lstlisting}

\begin{explanation}
\texttt{base} tells Vite that your site will be hosted at a subpath (username.github.io/repo-name/ instead of the root).
\end{explanation}

Install gh-pages package:

\begin{lstlisting}[language=bash]
npm install --save-dev gh-pages
\end{lstlisting}

Add deploy script to \texttt{package.json}:

\begin{lstlisting}[language=JavaScript,caption={Add to package.json scripts}]
"scripts": {
  "dev": "vite",
  "build": "tsc && vite build",
  "preview": "vite preview",
  "deploy": "npm run build && gh-pages -d dist"
}
\end{lstlisting}

Deploy:

\begin{lstlisting}[language=bash]
npm run deploy
\end{lstlisting}

\begin{checkpoint}
\begin{itemize}
    \item[$\square$] Wait a few minutes for GitHub to process
    \item[$\square$] Go to your repository Settings → Pages
    \item[$\square$] Ensure source is set to ``gh-pages'' branch
    \item[$\square$] Visit: https://yourusername.github.io/my-portfolio/
\end{itemize}
\end{checkpoint}

\begin{gitcommit}
\begin{lstlisting}[language=bash]
git add .
git commit -m "Configure deployment to GitHub Pages"
git push
\end{lstlisting}
\end{gitcommit}

%==============================================================================
% APPENDIX
%==============================================================================
\appendix

\section{Complete File Reference}

This appendix lists all files you should have created.

\subsection{File Tree}

\begin{lstlisting}[language=bash]
my-portfolio/
|-- public/
|-- src/
|   |-- assets/
|   |   |-- icons/
|   |-- components/
|   |   |-- Header.tsx
|   |   |-- Section.tsx
|   |   |-- ExperienceCard.tsx
|   |   |-- ResearchCard.tsx
|   |   |-- ProjectCard.tsx
|   |   |-- SocialLinks.tsx
|   |   |-- ThemeToggle.tsx
|   |   |-- SkillsGraph.tsx
|   |-- context/
|   |   |-- ThemeContext.tsx
|   |-- data/
|   |   |-- content.ts
|   |-- styles/
|   |   |-- Header.css
|   |   |-- Section.css
|   |   |-- ExperienceCard.css
|   |   |-- ResearchCard.css
|   |   |-- ProjectCard.css
|   |   |-- SocialLinks.css
|   |   |-- ThemeToggle.css
|   |   |-- SkillsGraph.css
|   |-- App.tsx
|   |-- App.css
|   |-- index.css
|   |-- main.tsx
|-- index.html
|-- package.json
|-- tsconfig.json
|-- vite.config.ts
\end{lstlisting}

\section{Git Commands Reference}

\begin{tabular}{|l|l|}
\hline
\textbf{Command} & \textbf{Purpose} \\
\hline
\texttt{git init} & Initialize repository \\
\texttt{git add .} & Stage all changes \\
\texttt{git commit -m "message"} & Create commit \\
\texttt{git push} & Upload to remote \\
\texttt{git status} & Check current state \\
\texttt{git log --oneline} & View commit history \\
\texttt{git diff} & See uncommitted changes \\
\hline
\end{tabular}

\section{Debugging Tips}

\subsection{Common Errors}

\begin{enumerate}
    \item \textbf{``Module not found''}: Check import path, ensure file exists
    \item \textbf{``X is not defined''}: Variable not declared or wrong scope
    \item \textbf{TypeScript errors}: Check type annotations match actual data
    \item \textbf{CSS not applying}: Check class name spelling, ensure CSS imported
    \item \textbf{Blank page}: Check browser console (F12) for errors
\end{enumerate}

\subsection{Using Browser DevTools}

\begin{enumerate}
    \item Press F12 to open DevTools
    \item \textbf{Console tab}: Shows errors and logs
    \item \textbf{Elements tab}: Inspect HTML structure and CSS
    \item \textbf{Network tab}: See what files are loading
    \item \textbf{React DevTools}: Install extension to inspect components
\end{enumerate}

\section{Next Steps}

After completing this tutorial, consider:

\begin{enumerate}
    \item Add animations with Framer Motion
    \item Add a blog section with Markdown support
    \item Improve accessibility (WCAG guidelines)
    \item Add SEO meta tags
    \item Set up a custom domain
    \item Add analytics (privacy-respecting like Plausible)
\end{enumerate}

\vspace{2cm}
\begin{center}
\textit{Congratulations on building your portfolio from scratch!}
\end{center}

\end{document}
